\documentclass[14pt]{extarticle}
\usepackage{fontspec}
\usepackage{polyglossia}
\usepackage{amsmath}
\usepackage{amsfonts}
\usepackage{amssymb}
\usepackage{pgfplots}
\pgfplotsset{compat=1.18}

% Hebrew setup
\setdefaultlanguage{hebrew}
\setotherlanguage{english}
\newfontfamily\hebrewfont[Script=Hebrew]{Arial}

% Math font setup for proper Greek symbols
\usepackage{unicode-math}
\setmathfont{Latin Modern Math}

% Increase line spacing throughout document
\linespread{1.3}

\usepackage[margin=1in]{geometry}

% Left-aligned equation environment
\newenvironment{lefteqs}{%
    \begin{english}%
    \begin{flushleft}%
    \renewcommand{\arraystretch}{1.3}%
    $\begin{array}{@{}l@{}}%
}{%
    \end{array}$%
    \end{flushleft}%
    \end{english}%
}

\begin{document}

\begin{center}
    {\LARGE \textbf{מטלה 2 - אלקטרומגנטיות אנליטית}}\\[0.5em]
    {\large שי רואימי - 318986270}\\[0.3em]
\end{center}

\hrule
\vspace{2em}

\textbf{\large שאלה 1}
(פרק 1 שאלה 2)

נתון הפוטנציאל $V=\frac{V_0}{b} (x^2 + y^2 + z^2)$

\textbf{א)} 
נמצא את גרדיינט הפוטנציאל בקוארדינטות קרטזיות:

\begin{lefteqs}
    \nabla V = \frac{V_0}{b} (2x \hat{x} + 2y \hat{y} + 2z \hat{z})
\end{lefteqs}

\textbf{ב)}
נראה כי הווקטור הנורמלי למשטחים שווי פוטנציאל שמצאתי הוא (בקאורדינטות כדוריות) ורטור היחידה בכיוון $r$

בקאורדינטות כדוריות, $V=\frac{V_0}{b} r^2$ (כי $r=\sqrt{x^2+y^2+z^2}$)

ולכן $\nabla V = 2\frac{V_0}{b} r \hat{r}$

לפי הפוטנציאל בקאורדינטות כדוריות, הוא גדל לפי $r^2$ והמשטחים שווי הפוטנציאל הם כדורים,
לכן הוקטור שמצביע לכיוון בו הפוטנציאל גדל הכי מהר הוא נורמלי למשטחים שווי הפוטנציאל ושווה ל $\nabla V$.
 ולכן הנורמל הוא:

 \begin{lefteqs}
    \frac{\nabla V}{|\nabla V|} = \frac{2\frac{V_0}{b} r \hat{r}}{2\frac{V_0}{b} r} = \hat{r}
 \end{lefteqs}

\textbf{ג)}
נבצע במפורש את אינטגרל המסלול על $\vec{E}$ בין ראשית הצירים לנק' $(a,a,0)$
לאורך המסילה $x=y \, , z=0$

נראה כי תוצאות האינטגרל אכן שוות להפרש הפוטנציאלים בין שתי הנקודות:

\begin{lefteqs}
    \vec{E} = -\nabla V(x,y,z) = -\frac{V_0}{b} (2x \hat{x} + 2y \hat{y} + 2z \hat{z}) \\
    x=y \, , z=0 \Rightarrow \vec{E} = -\frac{2V_0}{b} (x \hat{x} + x \hat{y}) \\
    d\vec{l} = dx\hat{x} + dy\hat{y} \\
    \int_{(0,0,0)}^{(a,a,0)} \vec{E} \cdot d\vec{l} = \int_{0}^{a} -\frac{2V_0}{b} (x \hat{x} + x \hat{y}) \cdot (dx\hat{x} + dy\hat{y}) \\
    = \frac{-2V_0}{b} \int_{0}^{a} xdx + xdx = \frac{-4V_0}{b} \int_{0}^{a} xdx \\
    = \frac{-4V_0}{b} \left[ \frac{x^2}{2} \right]_{0}^{a} = \frac{-4V_0}{b} \cdot \frac{a^2}{2} = \frac{-2V_0 a^2}{b} \\
    V(0,0,0) - V(a,a,0) = 0 - \frac{V_0}{b} (a^2 + a^2) = \frac{-2V_0 a^2}{b} \\
    \Rightarrow \int_{(0,0,0)}^{(a,a,0)} \vec{E} \cdot d\vec{l} = \Delta V \, \, \blacksquare
\end{lefteqs}


\end{document}